\begin{lexlean}{RustLower}
\useglossary{lexlean.std.nat@1.2.0}
\importmodule{CalculusCheck}
\title{Natural number addition}
\begin{semanticmodule}
inductive{
  name: Clos,
  type_parameters: [],
  parameters: [],
  constructors: [
    {name: ClNil, fields: []},
    {name: ClCons, fields: [CalculusCheck::Ty, Nat, CalculusCheck::Ty, Clos]}
  ]
}

inductive{
  name: Edges,
  type_parameters: [],
  parameters: [],
  constructors: [
    {name: EdgeNil, fields: []},
    {name: EdgeCons, fields: [CalculusCheck::Ty, CalculusCheck::Ty, Edges]}
  ]
}

inductive{
  name: Facts,
  type_parameters: [],
  parameters: [],
  constructors: [
    {name: FactsNil, fields: []},
    {name: FactsCons, fields: [Bool, (List Nat), CalculusCheck::Ty, Facts]}
  ]
}

structure{
  name: Collect,
  type_parameters: [],
  parameters: [],
  fields: [{name: known, type: CalculusCheck::Ty}, {name: closures, type: Clos}]
}

structure{
  name: Reach,
  type_parameters: [],
  parameters: [],
  fields: [
    {name: own, type: Bool},
    {name: calls, type: (List Nat)},
    {name: applies, type: CalculusCheck::Ty}
  ]
}

definition{
  name: tyMember,
  parameters: [chain : CalculusCheck::Ty, item : CalculusCheck::Ty],
  result: Bool,
  body: (match chain
    | CalculusCheck::Ty.TUnit => false
    | CalculusCheck::Ty.TBool => false
    | CalculusCheck::Ty.TNat => false
    | CalculusCheck::Ty.TInt => false
    | CalculusCheck::Ty.TFixed d0 => false
    | CalculusCheck::Ty.TString => false
    | CalculusCheck::Ty.TBytes => false
    | CalculusCheck::Ty.TOrdering => false
    | CalculusCheck::Ty.TOption d0 => false
    | CalculusCheck::Ty.TResult d0 d1 => false
    | CalculusCheck::Ty.TList d0 => false
    | CalculusCheck::Ty.TPair d0 d1 => false
    | CalculusCheck::Ty.TAdt d0 => false
    | CalculusCheck::Ty.TFn d0 d1 => false
    | CalculusCheck::Ty.TNil => false
    | CalculusCheck::Ty.TCons head tail => (if (CalculusCheck::tyEq head item) then true else (tyMember tail item))),
  executable: true,
  recursive_argument: chain
}

definition{
  name: tyAppend,
  parameters: [left : CalculusCheck::Ty, right : CalculusCheck::Ty],
  result: CalculusCheck::Ty,
  body: (match left
    | CalculusCheck::Ty.TUnit => right
    | CalculusCheck::Ty.TBool => right
    | CalculusCheck::Ty.TNat => right
    | CalculusCheck::Ty.TInt => right
    | CalculusCheck::Ty.TFixed d0 => right
    | CalculusCheck::Ty.TString => right
    | CalculusCheck::Ty.TBytes => right
    | CalculusCheck::Ty.TOrdering => right
    | CalculusCheck::Ty.TOption d0 => right
    | CalculusCheck::Ty.TResult d0 d1 => right
    | CalculusCheck::Ty.TList d0 => right
    | CalculusCheck::Ty.TPair d0 d1 => right
    | CalculusCheck::Ty.TAdt d0 => right
    | CalculusCheck::Ty.TFn d0 d1 => right
    | CalculusCheck::Ty.TNil => right
    | CalculusCheck::Ty.TCons head tail => (ctor CalculusCheck::Ty.TCons head (tyAppend tail right))),
  executable: true,
  recursive_argument: left
}

definition{
  name: tyAddNew,
  parameters: [chain : CalculusCheck::Ty, item : CalculusCheck::Ty],
  result: CalculusCheck::Ty,
  body: (if (tyMember chain item) then chain else (ctor CalculusCheck::Ty.TCons item chain)),
  executable: true
}

definition{
  name: tyDrop,
  parameters: [chain : CalculusCheck::Ty, count : Nat],
  result: CalculusCheck::Ty,
  body: (match chain
    | CalculusCheck::Ty.TUnit => (ctor CalculusCheck::Ty.TNil)
    | CalculusCheck::Ty.TBool => (ctor CalculusCheck::Ty.TNil)
    | CalculusCheck::Ty.TNat => (ctor CalculusCheck::Ty.TNil)
    | CalculusCheck::Ty.TInt => (ctor CalculusCheck::Ty.TNil)
    | CalculusCheck::Ty.TFixed d0 => (ctor CalculusCheck::Ty.TNil)
    | CalculusCheck::Ty.TString => (ctor CalculusCheck::Ty.TNil)
    | CalculusCheck::Ty.TBytes => (ctor CalculusCheck::Ty.TNil)
    | CalculusCheck::Ty.TOrdering => (ctor CalculusCheck::Ty.TNil)
    | CalculusCheck::Ty.TOption d0 => (ctor CalculusCheck::Ty.TNil)
    | CalculusCheck::Ty.TResult d0 d1 => (ctor CalculusCheck::Ty.TNil)
    | CalculusCheck::Ty.TList d0 => (ctor CalculusCheck::Ty.TNil)
    | CalculusCheck::Ty.TPair d0 d1 => (ctor CalculusCheck::Ty.TNil)
    | CalculusCheck::Ty.TAdt d0 => (ctor CalculusCheck::Ty.TNil)
    | CalculusCheck::Ty.TFn d0 d1 => (ctor CalculusCheck::Ty.TNil)
    | CalculusCheck::Ty.TNil => (ctor CalculusCheck::Ty.TNil)
    | CalculusCheck::Ty.TCons head tail => (match count
      | Nat.zero => (ctor CalculusCheck::Ty.TCons head tail)
      | Nat.succ previous => (tyDrop tail previous))),
  executable: true,
  recursive_argument: chain
}

definition{
  name: tyTake,
  parameters: [chain : CalculusCheck::Ty, count : Nat],
  result: CalculusCheck::Ty,
  body: (match chain
    | CalculusCheck::Ty.TUnit => (ctor CalculusCheck::Ty.TNil)
    | CalculusCheck::Ty.TBool => (ctor CalculusCheck::Ty.TNil)
    | CalculusCheck::Ty.TNat => (ctor CalculusCheck::Ty.TNil)
    | CalculusCheck::Ty.TInt => (ctor CalculusCheck::Ty.TNil)
    | CalculusCheck::Ty.TFixed d0 => (ctor CalculusCheck::Ty.TNil)
    | CalculusCheck::Ty.TString => (ctor CalculusCheck::Ty.TNil)
    | CalculusCheck::Ty.TBytes => (ctor CalculusCheck::Ty.TNil)
    | CalculusCheck::Ty.TOrdering => (ctor CalculusCheck::Ty.TNil)
    | CalculusCheck::Ty.TOption d0 => (ctor CalculusCheck::Ty.TNil)
    | CalculusCheck::Ty.TResult d0 d1 => (ctor CalculusCheck::Ty.TNil)
    | CalculusCheck::Ty.TList d0 => (ctor CalculusCheck::Ty.TNil)
    | CalculusCheck::Ty.TPair d0 d1 => (ctor CalculusCheck::Ty.TNil)
    | CalculusCheck::Ty.TAdt d0 => (ctor CalculusCheck::Ty.TNil)
    | CalculusCheck::Ty.TFn d0 d1 => (ctor CalculusCheck::Ty.TNil)
    | CalculusCheck::Ty.TNil => (ctor CalculusCheck::Ty.TNil)
    | CalculusCheck::Ty.TCons head tail => (match count
      | Nat.zero => (ctor CalculusCheck::Ty.TNil)
      | Nat.succ previous => (ctor CalculusCheck::Ty.TCons head (tyTake tail previous)))),
  executable: true,
  recursive_argument: chain
}

definition{
  name: register,
  parameters: [known : CalculusCheck::Ty, ty : CalculusCheck::Ty],
  result: CalculusCheck::Ty,
  body: (match ty
    | CalculusCheck::Ty.TUnit => known
    | CalculusCheck::Ty.TBool => known
    | CalculusCheck::Ty.TNat => known
    | CalculusCheck::Ty.TInt => known
    | CalculusCheck::Ty.TFixed d0 => known
    | CalculusCheck::Ty.TString => known
    | CalculusCheck::Ty.TBytes => known
    | CalculusCheck::Ty.TOrdering => known
    | CalculusCheck::Ty.TOption value => (register known value)
    | CalculusCheck::Ty.TResult ok error => (register (register known ok) error)
    | CalculusCheck::Ty.TList element => (register known element)
    | CalculusCheck::Ty.TPair left right => (register (register known left) right)
    | CalculusCheck::Ty.TAdt d0 => known
    | CalculusCheck::Ty.TFn parameters result => (let added : CalculusCheck::Ty := (tyAddNew known (ctor CalculusCheck::Ty.TFn parameters result)); (register (register added parameters) result))
    | CalculusCheck::Ty.TNil => known
    | CalculusCheck::Ty.TCons head tail => (register (register known head) tail)),
  executable: true,
  recursive_argument: ty
}

definition{
  name: primFails,
  parameters: [operation : CalculusCheck::Prim],
  result: Bool,
  body: (match operation
    | CalculusCheck::Prim.NatAdd => true
    | CalculusCheck::Prim.NatSub => false
    | CalculusCheck::Prim.NatMul => true
    | CalculusCheck::Prim.NatQuot => false
    | CalculusCheck::Prim.NatRem => false
    | CalculusCheck::Prim.NatEq => false
    | CalculusCheck::Prim.NatLe => false
    | CalculusCheck::Prim.NatLt => false
    | CalculusCheck::Prim.IntAdd => true
    | CalculusCheck::Prim.IntSub => true
    | CalculusCheck::Prim.IntMul => true
    | CalculusCheck::Prim.IntNeg => true
    | CalculusCheck::Prim.IntQuot => true
    | CalculusCheck::Prim.IntRem => false
    | CalculusCheck::Prim.CheckedAdd => false
    | CalculusCheck::Prim.CheckedSub => false
    | CalculusCheck::Prim.CheckedMul => false
    | CalculusCheck::Prim.CheckedNeg => false
    | CalculusCheck::Prim.CheckedQuot => false
    | CalculusCheck::Prim.BitAnd => false
    | CalculusCheck::Prim.BitOr => false
    | CalculusCheck::Prim.BitXor => false
    | CalculusCheck::Prim.BitNot => false
    | CalculusCheck::Prim.ShiftLeft => false
    | CalculusCheck::Prim.ShiftRight => false
    | CalculusCheck::Prim.Equal => false
    | CalculusCheck::Prim.BoolNot => false
    | CalculusCheck::Prim.BoolAnd => false
    | CalculusCheck::Prim.BoolOr => false
    | CalculusCheck::Prim.Append => false
    | CalculusCheck::Prim.Length => false
    | CalculusCheck::Prim.Index => false
    | CalculusCheck::Prim.Slice => false
    | CalculusCheck::Prim.Utf8Encode => false
    | CalculusCheck::Prim.Utf8Decode => false
    | CalculusCheck::Prim.CompareBytes => false
    | CalculusCheck::Prim.SplitExact => false
    | CalculusCheck::Prim.Join => false
    | CalculusCheck::Prim.FormatDecimal => false
    | CalculusCheck::Prim.Compare => false
    | CalculusCheck::Prim.Convert d0 => false
    | CalculusCheck::Prim.ParseDecimal target => ((CalculusCheck::tyTag target) =? 3)),
  executable: true
}

definition{
  name: closHas,
  parameters: [closures : Clos, fnTy : CalculusCheck::Ty, function : Nat],
  result: Bool,
  body: (match closures
    | Clos.ClNil => false
    | Clos.ClCons known member captured rest => (if (member =? function) then (if (CalculusCheck::tyEq known fnTy) then true else (closHas rest fnTy function)) else (closHas rest fnTy function))),
  executable: true,
  recursive_argument: closures
}

definition{
  name: closAdd,
  parameters: [
    closures : Clos,
    fnTy : CalculusCheck::Ty,
    function : Nat,
    captured : CalculusCheck::Ty
  ],
  result: Clos,
  body: (if (closHas closures fnTy function) then closures else (ctor Clos.ClCons fnTy function captured closures)),
  executable: true
}

definition{
  name: closVariants,
  parameters: [closures : Clos, fnTy : CalculusCheck::Ty],
  result: CalculusCheck::Ty,
  body: (match closures
    | Clos.ClNil => (ctor CalculusCheck::Ty.TNil)
    | Clos.ClCons known member captured rest => (if (CalculusCheck::tyEq known fnTy) then (ctor CalculusCheck::Ty.TCons captured (closVariants rest fnTy)) else (closVariants rest fnTy))),
  executable: true,
  recursive_argument: closures
}

definition{
  name: siteOwn,
  parameters: [sites : CalculusCheck::Ty],
  result: CalculusCheck::Ty,
  body: (CalculusCheck::tyFirst sites),
  executable: true
}

definition{
  name: kidA,
  parameters: [sites : CalculusCheck::Ty],
  result: CalculusCheck::Ty,
  body: (CalculusCheck::tyFirst (CalculusCheck::tySecond sites)),
  executable: true
}

definition{
  name: kidB,
  parameters: [sites : CalculusCheck::Ty],
  result: CalculusCheck::Ty,
  body: (CalculusCheck::tyFirst (CalculusCheck::tySecond (CalculusCheck::tySecond sites))),
  executable: true
}

definition{
  name: kidC,
  parameters: [sites : CalculusCheck::Ty],
  result: CalculusCheck::Ty,
  body: (CalculusCheck::tyFirst (CalculusCheck::tySecond (CalculusCheck::tySecond (CalculusCheck::tySecond sites)))),
  executable: true
}

definition{
  name: collectTerm,
  parameters: [
    program : CalculusCheck::Program,
    state : Collect,
    expr : CalculusCheck::Expr,
    sites : CalculusCheck::Ty
  ],
  result: Collect,
  body: (match expr
    | CalculusCheck::Expr.EValue ty value => record{
      type: Collect,
      fields: [
        {field: known, value: (register (project state known) ty)},
        {field: closures, value: (project state closures)}
      ],
      type_arguments: []
    }
    | CalculusCheck::Expr.EVar name => state
    | CalculusCheck::Expr.ELet name ty bound body => (let registered : Collect := record{
      type: Collect,
      fields: [
        {field: known, value: (register (project state known) ty)},
        {field: closures, value: (project state closures)}
      ],
      type_arguments: []
    }; (let inner : Collect := (collectTerm program registered bound (kidA sites)); (collectTerm program inner body (kidB sites))))
    | CalculusCheck::Expr.ECond condition thenBranch elseBranch => (let one : Collect := (collectTerm program state condition (kidA sites)); (let two : Collect := (collectTerm program one thenBranch (kidB sites)); (collectTerm program two elseBranch (kidC sites))))
    | CalculusCheck::Expr.EMatch ty scrutinee arms => (let registered : Collect := record{
      type: Collect,
      fields: [
        {field: known, value: (register (project state known) ty)},
        {field: closures, value: (project state closures)}
      ],
      type_arguments: []
    }; (let inner : Collect := (collectTerm program registered scrutinee (kidA sites)); (collectTerm program inner arms (kidB sites))))
    | CalculusCheck::Expr.EBuild shape ty operands => (let registered : Collect := record{
      type: Collect,
      fields: [
        {field: known, value: (register (project state known) ty)},
        {field: closures, value: (project state closures)}
      ],
      type_arguments: []
    }; (collectTerm program registered operands (kidA sites)))
    | CalculusCheck::Expr.ECall function operands => (collectTerm program state operands (kidA sites))
    | CalculusCheck::Expr.EClosure function captures => (let inner : Collect := (collectTerm program state captures (kidA sites)); (let declared : CalculusCheck::Ty := (CalculusCheck::functionTypesOr program function); (let count : Nat := (CalculusCheck::exprChainLength captures); (let closureType : CalculusCheck::Ty := (ctor CalculusCheck::Ty.TFn (tyDrop declared count) (CalculusCheck::functionResult program function)); (let registered : Collect := record{
      type: Collect,
      fields: [
        {field: known, value: (register (project inner known) closureType)},
        {field: closures, value: (project inner closures)}
      ],
      type_arguments: []
    }; record{
      type: Collect,
      fields: [
        {field: known, value: (project registered known)},
        {
          field: closures,
          value: (closAdd (project registered closures) closureType function (tyTake declared count))
        }
      ],
      type_arguments: []
    })))))
    | CalculusCheck::Expr.EApply target operands => (let inner : Collect := (collectTerm program state target (kidA sites)); (collectTerm program inner operands (kidB sites)))
    | CalculusCheck::Expr.EPrim operation operands => (collectTerm program state operands (kidA sites))
    | CalculusCheck::Expr.EFirst value => (collectTerm program state value (kidA sites))
    | CalculusCheck::Expr.ESecond value => (collectTerm program state value (kidA sites))
    | CalculusCheck::Expr.EField value index => (collectTerm program state value (kidA sites))
    | CalculusCheck::Expr.ENil => state
    | CalculusCheck::Expr.ECons head tail => (let inner : Collect := (collectTerm program state head (kidA sites)); (collectTerm program inner tail (kidB sites)))
    | CalculusCheck::Expr.EArm shape binders body => (collectTerm program state body (kidA sites))),
  executable: true,
  recursive_argument: expr
}

definition{
  name: collectFunctions,
  parameters: [
    program : CalculusCheck::Program,
    state : Collect,
    bodies : CalculusCheck::Expr,
    index : Nat,
    dead : (List Bool),
    sites : CalculusCheck::Ty
  ],
  result: Collect,
  body: (match bodies
    | CalculusCheck::Expr.EValue d0 d1 => state
    | CalculusCheck::Expr.EVar d0 => state
    | CalculusCheck::Expr.ELet d0 d1 d2 d3 => state
    | CalculusCheck::Expr.ECond d0 d1 d2 => state
    | CalculusCheck::Expr.EMatch d0 d1 d2 => state
    | CalculusCheck::Expr.EBuild d0 d1 d2 => state
    | CalculusCheck::Expr.ECall d0 d1 => state
    | CalculusCheck::Expr.EClosure d0 d1 => state
    | CalculusCheck::Expr.EApply d0 d1 => state
    | CalculusCheck::Expr.EPrim d0 d1 => state
    | CalculusCheck::Expr.EFirst d0 => state
    | CalculusCheck::Expr.ESecond d0 => state
    | CalculusCheck::Expr.EField d0 d1 => state
    | CalculusCheck::Expr.ENil => state
    | CalculusCheck::Expr.ECons body rest => (let live : Bool := (match (CalculusCheck::listAt [Bool] dead index)
      | Option.some flag => (not flag)
      | Option.none => true); (let inner : Collect := (if live then (collectTerm program state body (kidA sites)) else state); (collectFunctions program inner rest (index + 1) dead (kidB sites))))
    | CalculusCheck::Expr.EArm d0 d1 d2 => state),
  executable: true,
  recursive_argument: bodies
}

definition{
  name: natMember,
  parameters: [items : (List Nat), item : Nat],
  result: Bool,
  body: (match items
    | List.nil => false
    | List.cons head tail => (if (head =? item) then true else (natMember tail item))),
  executable: true,
  recursive_argument: items
}

definition{
  name: reachTerm,
  parameters: [
    program : CalculusCheck::Program,
    found : Reach,
    expr : CalculusCheck::Expr,
    sites : CalculusCheck::Ty
  ],
  result: Reach,
  body: (match expr
    | CalculusCheck::Expr.EValue ty value => found
    | CalculusCheck::Expr.EVar name => found
    | CalculusCheck::Expr.ELet name ty bound body => (let inner : Reach := (reachTerm program found bound (kidA sites)); (reachTerm program inner body (kidB sites)))
    | CalculusCheck::Expr.ECond condition thenBranch elseBranch => (let one : Reach := (reachTerm program found condition (kidA sites)); (let two : Reach := (reachTerm program one thenBranch (kidB sites)); (reachTerm program two elseBranch (kidC sites))))
    | CalculusCheck::Expr.EMatch ty scrutinee arms => (let inner : Reach := (reachTerm program found scrutinee (kidA sites)); (reachTerm program inner arms (kidB sites)))
    | CalculusCheck::Expr.EBuild shape ty operands => (let inner : Reach := (reachTerm program found operands (kidA sites)); (match shape
      | CalculusCheck::Shape.SSucc => record{
        type: Reach,
        fields: [
          {field: own, value: true},
          {field: calls, value: (project inner calls)},
          {field: applies, value: (project inner applies)}
        ],
        type_arguments: []
      }
      | CalculusCheck::Shape.SNone => inner
      | CalculusCheck::Shape.SSome => inner
      | CalculusCheck::Shape.SOk => inner
      | CalculusCheck::Shape.SError => inner
      | CalculusCheck::Shape.SNil => inner
      | CalculusCheck::Shape.SCons => inner
      | CalculusCheck::Shape.SZero => inner
      | CalculusCheck::Shape.SPair => inner
      | CalculusCheck::Shape.STrue => inner
      | CalculusCheck::Shape.SFalse => inner
      | CalculusCheck::Shape.SUnit => inner
      | CalculusCheck::Shape.SLt => inner
      | CalculusCheck::Shape.SEq => inner
      | CalculusCheck::Shape.SGt => inner
      | CalculusCheck::Shape.SAdt d0 => inner))
    | CalculusCheck::Expr.ECall function operands => (let inner : Reach := (reachTerm program found operands (kidA sites)); record{
      type: Reach,
      fields: [
        {field: own, value: (project inner own)},
        {
          field: calls,
          value: (if (natMember (project inner calls) function) then (project inner calls) else (cons function (project inner calls)))
        },
        {field: applies, value: (project inner applies)}
      ],
      type_arguments: []
    })
    | CalculusCheck::Expr.EClosure function captures => (reachTerm program found captures (kidA sites))
    | CalculusCheck::Expr.EApply target operands => (let one : Reach := (reachTerm program found target (kidA sites)); (let two : Reach := (reachTerm program one operands (kidB sites)); record{
      type: Reach,
      fields: [
        {field: own, value: (project two own)},
        {field: calls, value: (project two calls)},
        {field: applies, value: (tyAddNew (project two applies) (siteOwn sites))}
      ],
      type_arguments: []
    }))
    | CalculusCheck::Expr.EPrim operation operands => (let inner : Reach := (reachTerm program found operands (kidA sites)); record{
      type: Reach,
      fields: [
        {field: own, value: (if (project inner own) then true else (primFails operation))},
        {field: calls, value: (project inner calls)},
        {field: applies, value: (project inner applies)}
      ],
      type_arguments: []
    })
    | CalculusCheck::Expr.EFirst value => (reachTerm program found value (kidA sites))
    | CalculusCheck::Expr.ESecond value => (reachTerm program found value (kidA sites))
    | CalculusCheck::Expr.EField value index => (reachTerm program found value (kidA sites))
    | CalculusCheck::Expr.ENil => found
    | CalculusCheck::Expr.ECons head tail => (let inner : Reach := (reachTerm program found head (kidA sites)); (reachTerm program inner tail (kidB sites)))
    | CalculusCheck::Expr.EArm shape binders body => (reachTerm program found body (kidA sites))),
  executable: true,
  recursive_argument: expr
}

definition{
  name: reachFunctions,
  parameters: [
    program : CalculusCheck::Program,
    bodies : CalculusCheck::Expr,
    index : Nat,
    sites : CalculusCheck::Ty
  ],
  result: Facts,
  body: (match bodies
    | CalculusCheck::Expr.EValue d0 d1 => (ctor Facts.FactsNil)
    | CalculusCheck::Expr.EVar d0 => (ctor Facts.FactsNil)
    | CalculusCheck::Expr.ELet d0 d1 d2 d3 => (ctor Facts.FactsNil)
    | CalculusCheck::Expr.ECond d0 d1 d2 => (ctor Facts.FactsNil)
    | CalculusCheck::Expr.EMatch d0 d1 d2 => (ctor Facts.FactsNil)
    | CalculusCheck::Expr.EBuild d0 d1 d2 => (ctor Facts.FactsNil)
    | CalculusCheck::Expr.ECall d0 d1 => (ctor Facts.FactsNil)
    | CalculusCheck::Expr.EClosure d0 d1 => (ctor Facts.FactsNil)
    | CalculusCheck::Expr.EApply d0 d1 => (ctor Facts.FactsNil)
    | CalculusCheck::Expr.EPrim d0 d1 => (ctor Facts.FactsNil)
    | CalculusCheck::Expr.EFirst d0 => (ctor Facts.FactsNil)
    | CalculusCheck::Expr.ESecond d0 => (ctor Facts.FactsNil)
    | CalculusCheck::Expr.EField d0 d1 => (ctor Facts.FactsNil)
    | CalculusCheck::Expr.ENil => (ctor Facts.FactsNil)
    | CalculusCheck::Expr.ECons body rest => (let found : Reach := (reachTerm program record{
      type: Reach,
      fields: [
        {field: own, value: false},
        {field: calls, value: (nil Nat)},
        {field: applies, value: (ctor CalculusCheck::Ty.TNil)}
      ],
      type_arguments: []
    } body (kidA sites)); (ctor Facts.FactsCons (project found own) (project found calls) (project found applies) (reachFunctions program rest (index + 1) (kidB sites))))
    | CalculusCheck::Expr.EArm d0 d1 d2 => (ctor Facts.FactsNil)),
  executable: true,
  recursive_argument: bodies
}

definition{
  name: adtOwners,
  parameters: [count : Nat],
  result: CalculusCheck::Ty,
  body: (match count
    | Nat.zero => (ctor CalculusCheck::Ty.TNil)
    | Nat.succ previous => (ctor CalculusCheck::Ty.TCons (ctor CalculusCheck::Ty.TAdt previous) (adtOwners previous))),
  executable: true,
  recursive_argument: count
}

definition{
  name: inPlace,
  parameters: [known : CalculusCheck::Ty, held : CalculusCheck::Ty, ty : CalculusCheck::Ty],
  result: CalculusCheck::Ty,
  body: (match ty
    | CalculusCheck::Ty.TUnit => held
    | CalculusCheck::Ty.TBool => held
    | CalculusCheck::Ty.TNat => held
    | CalculusCheck::Ty.TInt => held
    | CalculusCheck::Ty.TFixed d0 => held
    | CalculusCheck::Ty.TString => held
    | CalculusCheck::Ty.TBytes => held
    | CalculusCheck::Ty.TOrdering => held
    | CalculusCheck::Ty.TOption value => (inPlace known held value)
    | CalculusCheck::Ty.TResult ok error => (inPlace known (inPlace known held ok) error)
    | CalculusCheck::Ty.TList d0 => held
    | CalculusCheck::Ty.TPair left right => (inPlace known (inPlace known held left) right)
    | CalculusCheck::Ty.TAdt index => (tyAddNew held (ctor CalculusCheck::Ty.TAdt index))
    | CalculusCheck::Ty.TFn parameters result => (if (tyMember known (ctor CalculusCheck::Ty.TFn parameters result)) then (tyAddNew held (ctor CalculusCheck::Ty.TFn parameters result)) else held)
    | CalculusCheck::Ty.TNil => held
    | CalculusCheck::Ty.TCons d0 d1 => held),
  executable: true,
  recursive_argument: ty
}

definition{
  name: edgeHas,
  parameters: [edges : Edges, from : CalculusCheck::Ty, to : CalculusCheck::Ty],
  result: Bool,
  body: (match edges
    | Edges.EdgeNil => false
    | Edges.EdgeCons left right rest => (if (CalculusCheck::tyEq left from) then (if (CalculusCheck::tyEq right to) then true else (edgeHas rest from to)) else (edgeHas rest from to))),
  executable: true,
  recursive_argument: edges
}

definition{
  name: edgeCount,
  parameters: [edges : Edges],
  result: Nat,
  body: (match edges
    | Edges.EdgeNil => 0
    | Edges.EdgeCons left right rest => (1 + (edgeCount rest))),
  executable: true,
  recursive_argument: edges
}

definition{
  name: edgeAdd,
  parameters: [edges : Edges, from : CalculusCheck::Ty, to : CalculusCheck::Ty],
  result: Edges,
  body: (if (edgeHas edges from to) then edges else (ctor Edges.EdgeCons from to edges)),
  executable: true
}

definition{
  name: edgesToHeld,
  parameters: [owner : CalculusCheck::Ty, held : CalculusCheck::Ty, edges : Edges],
  result: Edges,
  body: (match held
    | CalculusCheck::Ty.TUnit => edges
    | CalculusCheck::Ty.TBool => edges
    | CalculusCheck::Ty.TNat => edges
    | CalculusCheck::Ty.TInt => edges
    | CalculusCheck::Ty.TFixed d0 => edges
    | CalculusCheck::Ty.TString => edges
    | CalculusCheck::Ty.TBytes => edges
    | CalculusCheck::Ty.TOrdering => edges
    | CalculusCheck::Ty.TOption d0 => edges
    | CalculusCheck::Ty.TResult d0 d1 => edges
    | CalculusCheck::Ty.TList d0 => edges
    | CalculusCheck::Ty.TPair d0 d1 => edges
    | CalculusCheck::Ty.TAdt d0 => edges
    | CalculusCheck::Ty.TFn d0 d1 => edges
    | CalculusCheck::Ty.TNil => edges
    | CalculusCheck::Ty.TCons head tail => (edgesToHeld owner tail (edgeAdd edges owner head))),
  executable: true,
  recursive_argument: held
}

definition{
  name: edgesOfFields,
  parameters: [
    known : CalculusCheck::Ty,
    owner : CalculusCheck::Ty,
    fields : CalculusCheck::Ty,
    edges : Edges
  ],
  result: Edges,
  body: (match fields
    | CalculusCheck::Ty.TUnit => edges
    | CalculusCheck::Ty.TBool => edges
    | CalculusCheck::Ty.TNat => edges
    | CalculusCheck::Ty.TInt => edges
    | CalculusCheck::Ty.TFixed d0 => edges
    | CalculusCheck::Ty.TString => edges
    | CalculusCheck::Ty.TBytes => edges
    | CalculusCheck::Ty.TOrdering => edges
    | CalculusCheck::Ty.TOption d0 => edges
    | CalculusCheck::Ty.TResult d0 d1 => edges
    | CalculusCheck::Ty.TList d0 => edges
    | CalculusCheck::Ty.TPair d0 d1 => edges
    | CalculusCheck::Ty.TAdt d0 => edges
    | CalculusCheck::Ty.TFn d0 d1 => edges
    | CalculusCheck::Ty.TNil => edges
    | CalculusCheck::Ty.TCons head tail => (edgesOfFields known owner tail (edgesToHeld owner (inPlace known (ctor CalculusCheck::Ty.TNil) head) edges))),
  executable: true,
  recursive_argument: fields
}

definition{
  name: edgesOfVariants,
  parameters: [
    known : CalculusCheck::Ty,
    owner : CalculusCheck::Ty,
    variants : CalculusCheck::Ty,
    edges : Edges
  ],
  result: Edges,
  body: (match variants
    | CalculusCheck::Ty.TUnit => edges
    | CalculusCheck::Ty.TBool => edges
    | CalculusCheck::Ty.TNat => edges
    | CalculusCheck::Ty.TInt => edges
    | CalculusCheck::Ty.TFixed d0 => edges
    | CalculusCheck::Ty.TString => edges
    | CalculusCheck::Ty.TBytes => edges
    | CalculusCheck::Ty.TOrdering => edges
    | CalculusCheck::Ty.TOption d0 => edges
    | CalculusCheck::Ty.TResult d0 d1 => edges
    | CalculusCheck::Ty.TList d0 => edges
    | CalculusCheck::Ty.TPair d0 d1 => edges
    | CalculusCheck::Ty.TAdt d0 => edges
    | CalculusCheck::Ty.TFn d0 d1 => edges
    | CalculusCheck::Ty.TNil => edges
    | CalculusCheck::Ty.TCons head tail => (edgesOfVariants known owner tail (edgesOfFields known owner head edges))),
  executable: true,
  recursive_argument: variants
}

definition{
  name: variantsOf,
  parameters: [program : CalculusCheck::Program, closures : Clos, owner : CalculusCheck::Ty],
  result: CalculusCheck::Ty,
  body: (match owner
    | CalculusCheck::Ty.TUnit => (ctor CalculusCheck::Ty.TNil)
    | CalculusCheck::Ty.TBool => (ctor CalculusCheck::Ty.TNil)
    | CalculusCheck::Ty.TNat => (ctor CalculusCheck::Ty.TNil)
    | CalculusCheck::Ty.TInt => (ctor CalculusCheck::Ty.TNil)
    | CalculusCheck::Ty.TFixed d0 => (ctor CalculusCheck::Ty.TNil)
    | CalculusCheck::Ty.TString => (ctor CalculusCheck::Ty.TNil)
    | CalculusCheck::Ty.TBytes => (ctor CalculusCheck::Ty.TNil)
    | CalculusCheck::Ty.TOrdering => (ctor CalculusCheck::Ty.TNil)
    | CalculusCheck::Ty.TOption d0 => (ctor CalculusCheck::Ty.TNil)
    | CalculusCheck::Ty.TResult d0 d1 => (ctor CalculusCheck::Ty.TNil)
    | CalculusCheck::Ty.TList d0 => (ctor CalculusCheck::Ty.TNil)
    | CalculusCheck::Ty.TPair d0 d1 => (ctor CalculusCheck::Ty.TNil)
    | CalculusCheck::Ty.TAdt index => (match (CalculusCheck::adtAt program index)
      | Option.some constructors => constructors
      | Option.none => (ctor CalculusCheck::Ty.TNil))
    | CalculusCheck::Ty.TFn parameters result => (closVariants closures (ctor CalculusCheck::Ty.TFn parameters result))
    | CalculusCheck::Ty.TNil => (ctor CalculusCheck::Ty.TNil)
    | CalculusCheck::Ty.TCons d0 d1 => (ctor CalculusCheck::Ty.TNil)),
  executable: true
}

definition{
  name: allEdges,
  parameters: [
    program : CalculusCheck::Program,
    known : CalculusCheck::Ty,
    closures : Clos,
    owners : CalculusCheck::Ty,
    edges : Edges
  ],
  result: Edges,
  body: (match owners
    | CalculusCheck::Ty.TUnit => edges
    | CalculusCheck::Ty.TBool => edges
    | CalculusCheck::Ty.TNat => edges
    | CalculusCheck::Ty.TInt => edges
    | CalculusCheck::Ty.TFixed d0 => edges
    | CalculusCheck::Ty.TString => edges
    | CalculusCheck::Ty.TBytes => edges
    | CalculusCheck::Ty.TOrdering => edges
    | CalculusCheck::Ty.TOption d0 => edges
    | CalculusCheck::Ty.TResult d0 d1 => edges
    | CalculusCheck::Ty.TList d0 => edges
    | CalculusCheck::Ty.TPair d0 d1 => edges
    | CalculusCheck::Ty.TAdt d0 => edges
    | CalculusCheck::Ty.TFn d0 d1 => edges
    | CalculusCheck::Ty.TNil => edges
    | CalculusCheck::Ty.TCons head tail => (allEdges program known closures tail (edgesOfVariants known head (variantsOf program closures head) edges))),
  executable: true,
  recursive_argument: owners
}

definition{
  name: composeFrom,
  parameters: [right : Edges, middle : CalculusCheck::Ty, from : CalculusCheck::Ty, edges : Edges],
  result: Edges,
  body: (match right
    | Edges.EdgeNil => edges
    | Edges.EdgeCons start to rest => (composeFrom rest middle from (if (CalculusCheck::tyEq start middle) then (edgeAdd edges from to) else edges))),
  executable: true,
  recursive_argument: right
}

definition{
  name: compose,
  parameters: [all : Edges, left : Edges, edges : Edges],
  result: Edges,
  body: (match left
    | Edges.EdgeNil => edges
    | Edges.EdgeCons from middle rest => (compose all rest (composeFrom all middle from edges))),
  executable: true,
  recursive_argument: left
}

definition{
  name: closeStep,
  parameters: [edges : Edges],
  result: (Option Edges),
  body: (let wider : Edges := (compose edges edges edges); (if ((edgeCount wider) =? (edgeCount edges)) then (ctor Option.none [Edges]) else (ctor Option.some [Edges] wider))),
  executable: true
}

definition{
  name: closeEdges,
  parameters: [edges : Edges, bound : Nat],
  result: Edges,
  body: (first (prim iterate_until (fun (current : Edges) => (closeStep current)) bound edges : (Edges * Bool))),
  executable: true
}

definition{
  name: visible,
  parameters: [ty : CalculusCheck::Ty, inhabited : CalculusCheck::Ty],
  result: Bool,
  body: (match ty
    | CalculusCheck::Ty.TUnit => true
    | CalculusCheck::Ty.TBool => true
    | CalculusCheck::Ty.TNat => true
    | CalculusCheck::Ty.TInt => true
    | CalculusCheck::Ty.TFixed d0 => true
    | CalculusCheck::Ty.TString => true
    | CalculusCheck::Ty.TBytes => true
    | CalculusCheck::Ty.TOrdering => true
    | CalculusCheck::Ty.TOption d0 => true
    | CalculusCheck::Ty.TResult ok error => (if (visible ok inhabited) then true else (visible error inhabited))
    | CalculusCheck::Ty.TList d0 => true
    | CalculusCheck::Ty.TPair left right => (if (visible left inhabited) then (visible right inhabited) else false)
    | CalculusCheck::Ty.TAdt index => (tyMember inhabited (ctor CalculusCheck::Ty.TAdt index))
    | CalculusCheck::Ty.TFn parameters result => (tyMember inhabited (ctor CalculusCheck::Ty.TFn parameters result))
    | CalculusCheck::Ty.TNil => true
    | CalculusCheck::Ty.TCons d0 d1 => true),
  executable: true,
  recursive_argument: ty
}

definition{
  name: anyHeld,
  parameters: [owner : CalculusCheck::Ty, edges : Edges, held : CalculusCheck::Ty],
  result: Bool,
  body: (match held
    | CalculusCheck::Ty.TUnit => false
    | CalculusCheck::Ty.TBool => false
    | CalculusCheck::Ty.TNat => false
    | CalculusCheck::Ty.TInt => false
    | CalculusCheck::Ty.TFixed d0 => false
    | CalculusCheck::Ty.TString => false
    | CalculusCheck::Ty.TBytes => false
    | CalculusCheck::Ty.TOrdering => false
    | CalculusCheck::Ty.TOption d0 => false
    | CalculusCheck::Ty.TResult d0 d1 => false
    | CalculusCheck::Ty.TList d0 => false
    | CalculusCheck::Ty.TPair d0 d1 => false
    | CalculusCheck::Ty.TAdt d0 => false
    | CalculusCheck::Ty.TFn d0 d1 => false
    | CalculusCheck::Ty.TNil => false
    | CalculusCheck::Ty.TCons head tail => (if (CalculusCheck::tyEq head owner) then true else (if (edgeHas edges head owner) then true else (anyHeld owner edges tail)))),
  executable: true,
  recursive_argument: held
}

definition{
  name: cyclic,
  parameters: [
    known : CalculusCheck::Ty,
    edges : Edges,
    owner : CalculusCheck::Ty,
    ty : CalculusCheck::Ty
  ],
  result: Bool,
  body: (anyHeld owner edges (inPlace known (ctor CalculusCheck::Ty.TNil) ty)),
  executable: true
}

definition{
  name: fieldsOk,
  parameters: [
    known : CalculusCheck::Ty,
    edges : Edges,
    inhabited : CalculusCheck::Ty,
    owner : CalculusCheck::Ty,
    fields : CalculusCheck::Ty
  ],
  result: Bool,
  body: (match fields
    | CalculusCheck::Ty.TUnit => true
    | CalculusCheck::Ty.TBool => true
    | CalculusCheck::Ty.TNat => true
    | CalculusCheck::Ty.TInt => true
    | CalculusCheck::Ty.TFixed d0 => true
    | CalculusCheck::Ty.TString => true
    | CalculusCheck::Ty.TBytes => true
    | CalculusCheck::Ty.TOrdering => true
    | CalculusCheck::Ty.TOption d0 => true
    | CalculusCheck::Ty.TResult d0 d1 => true
    | CalculusCheck::Ty.TList d0 => true
    | CalculusCheck::Ty.TPair d0 d1 => true
    | CalculusCheck::Ty.TAdt d0 => true
    | CalculusCheck::Ty.TFn d0 d1 => true
    | CalculusCheck::Ty.TNil => true
    | CalculusCheck::Ty.TCons head tail => (if (if (cyclic known edges owner head) then true else (visible head inhabited)) then (fieldsOk known edges inhabited owner tail) else false)),
  executable: true,
  recursive_argument: fields
}

definition{
  name: someVariantOk,
  parameters: [
    known : CalculusCheck::Ty,
    edges : Edges,
    inhabited : CalculusCheck::Ty,
    owner : CalculusCheck::Ty,
    variants : CalculusCheck::Ty
  ],
  result: Bool,
  body: (match variants
    | CalculusCheck::Ty.TUnit => false
    | CalculusCheck::Ty.TBool => false
    | CalculusCheck::Ty.TNat => false
    | CalculusCheck::Ty.TInt => false
    | CalculusCheck::Ty.TFixed d0 => false
    | CalculusCheck::Ty.TString => false
    | CalculusCheck::Ty.TBytes => false
    | CalculusCheck::Ty.TOrdering => false
    | CalculusCheck::Ty.TOption d0 => false
    | CalculusCheck::Ty.TResult d0 d1 => false
    | CalculusCheck::Ty.TList d0 => false
    | CalculusCheck::Ty.TPair d0 d1 => false
    | CalculusCheck::Ty.TAdt d0 => false
    | CalculusCheck::Ty.TFn d0 d1 => false
    | CalculusCheck::Ty.TNil => false
    | CalculusCheck::Ty.TCons head tail => (if (fieldsOk known edges inhabited owner head) then true else (someVariantOk known edges inhabited owner tail))),
  executable: true,
  recursive_argument: variants
}

definition{
  name: inhabitRound,
  parameters: [
    program : CalculusCheck::Program,
    known : CalculusCheck::Ty,
    closures : Clos,
    edges : Edges,
    owners : CalculusCheck::Ty,
    inhabited : CalculusCheck::Ty
  ],
  result: CalculusCheck::Ty,
  body: (match owners
    | CalculusCheck::Ty.TUnit => inhabited
    | CalculusCheck::Ty.TBool => inhabited
    | CalculusCheck::Ty.TNat => inhabited
    | CalculusCheck::Ty.TInt => inhabited
    | CalculusCheck::Ty.TFixed d0 => inhabited
    | CalculusCheck::Ty.TString => inhabited
    | CalculusCheck::Ty.TBytes => inhabited
    | CalculusCheck::Ty.TOrdering => inhabited
    | CalculusCheck::Ty.TOption d0 => inhabited
    | CalculusCheck::Ty.TResult d0 d1 => inhabited
    | CalculusCheck::Ty.TList d0 => inhabited
    | CalculusCheck::Ty.TPair d0 d1 => inhabited
    | CalculusCheck::Ty.TAdt d0 => inhabited
    | CalculusCheck::Ty.TFn d0 d1 => inhabited
    | CalculusCheck::Ty.TNil => inhabited
    | CalculusCheck::Ty.TCons head tail => (inhabitRound program known closures edges tail (if (tyMember inhabited head) then inhabited else (if (someVariantOk known edges inhabited head (variantsOf program closures head)) then (tyAddNew inhabited head) else inhabited)))),
  executable: true,
  recursive_argument: owners
}

definition{
  name: inhabitStep,
  parameters: [
    program : CalculusCheck::Program,
    known : CalculusCheck::Ty,
    closures : Clos,
    edges : Edges,
    owners : CalculusCheck::Ty,
    inhabited : CalculusCheck::Ty
  ],
  result: (Option CalculusCheck::Ty),
  body: (let next : CalculusCheck::Ty := (inhabitRound program known closures edges owners inhabited); (if ((CalculusCheck::tyChainLength next) =? (CalculusCheck::tyChainLength inhabited)) then (ctor Option.none [CalculusCheck::Ty]) else (ctor Option.some [CalculusCheck::Ty] next))),
  executable: true
}

definition{
  name: inhabitation,
  parameters: [
    program : CalculusCheck::Program,
    known : CalculusCheck::Ty,
    closures : Clos,
    edges : Edges,
    owners : CalculusCheck::Ty
  ],
  result: CalculusCheck::Ty,
  body: (first (prim iterate_until (fun [closures, edges, known, owners, program] (inhabited : CalculusCheck::Ty) => (inhabitStep program known closures edges owners inhabited)) (1 + (CalculusCheck::tyChainLength owners)) (ctor CalculusCheck::Ty.TNil) : (CalculusCheck::Ty * Bool))),
  executable: true
}

definition{
  name: anyInvisible,
  parameters: [params : CalculusCheck::Ty, inhabited : CalculusCheck::Ty],
  result: Bool,
  body: (match params
    | CalculusCheck::Ty.TUnit => false
    | CalculusCheck::Ty.TBool => false
    | CalculusCheck::Ty.TNat => false
    | CalculusCheck::Ty.TInt => false
    | CalculusCheck::Ty.TFixed d0 => false
    | CalculusCheck::Ty.TString => false
    | CalculusCheck::Ty.TBytes => false
    | CalculusCheck::Ty.TOrdering => false
    | CalculusCheck::Ty.TOption d0 => false
    | CalculusCheck::Ty.TResult d0 d1 => false
    | CalculusCheck::Ty.TList d0 => false
    | CalculusCheck::Ty.TPair d0 d1 => false
    | CalculusCheck::Ty.TAdt d0 => false
    | CalculusCheck::Ty.TFn d0 d1 => false
    | CalculusCheck::Ty.TNil => false
    | CalculusCheck::Ty.TCons head tail => (if (visible head inhabited) then (anyInvisible tail inhabited) else true)),
  executable: true,
  recursive_argument: params
}

definition{
  name: deadFlags,
  parameters: [types : CalculusCheck::Ty, inhabited : CalculusCheck::Ty],
  result: (List Bool),
  body: (match types
    | CalculusCheck::Ty.TUnit => (nil Bool)
    | CalculusCheck::Ty.TBool => (nil Bool)
    | CalculusCheck::Ty.TNat => (nil Bool)
    | CalculusCheck::Ty.TInt => (nil Bool)
    | CalculusCheck::Ty.TFixed d0 => (nil Bool)
    | CalculusCheck::Ty.TString => (nil Bool)
    | CalculusCheck::Ty.TBytes => (nil Bool)
    | CalculusCheck::Ty.TOrdering => (nil Bool)
    | CalculusCheck::Ty.TOption d0 => (nil Bool)
    | CalculusCheck::Ty.TResult d0 d1 => (nil Bool)
    | CalculusCheck::Ty.TList d0 => (nil Bool)
    | CalculusCheck::Ty.TPair d0 d1 => (nil Bool)
    | CalculusCheck::Ty.TAdt d0 => (nil Bool)
    | CalculusCheck::Ty.TFn d0 d1 => (nil Bool)
    | CalculusCheck::Ty.TNil => (nil Bool)
    | CalculusCheck::Ty.TCons head tail => (cons (anyInvisible head inhabited) (deadFlags tail inhabited))),
  executable: true,
  recursive_argument: types
}

definition{
  name: boolsEq,
  parameters: [left : (List Bool), right : (List Bool)],
  result: Bool,
  body: (match left
    | List.nil => (match right | List.nil => true | List.cons rh rt => false)
    | List.cons lh lt => (match right
      | List.nil => false
      | List.cons rh rt => (if (prim equal lh rh : Bool) then (boolsEq lt rt) else false))),
  executable: true,
  recursive_argument: left
}

structure{
  name: Liveness,
  type_parameters: [],
  parameters: [],
  fields: [{name: dead, type: (List Bool)}, {name: closures, type: Clos}]
}

definition{
  name: livenessStep,
  parameters: [
    program : CalculusCheck::Program,
    bodies : CalculusCheck::Expr,
    known : CalculusCheck::Ty,
    edges : Edges,
    owners : CalculusCheck::Ty,
    sites : CalculusCheck::Ty,
    state : Liveness
  ],
  result: (Option Liveness),
  body: (let inhabited : CalculusCheck::Ty := (inhabitation program known (project state closures) edges owners); (let dead : (List Bool) := (deadFlags (project program types) inhabited); (if (boolsEq dead (project state dead)) then (ctor Option.none [Liveness]) else (ctor Option.some [Liveness] record{
    type: Liveness,
    fields: [
      {field: dead, value: dead},
      {
        field: closures,
        value: (project (collectFunctions program record{
          type: Collect,
          fields: [{field: known, value: known}, {field: closures, value: (ctor Clos.ClNil)}],
          type_arguments: []
        } bodies 0 dead sites) closures)
      }
    ],
    type_arguments: []
  })))),
  executable: true
}

definition{
  name: falses,
  parameters: [count : Nat],
  result: (List Bool),
  body: (match count
    | Nat.zero => (nil Bool)
    | Nat.succ previous => (cons false (falses previous))),
  executable: true,
  recursive_argument: count
}

definition{
  name: flagAt,
  parameters: [flags : (List Bool), index : Nat],
  result: Bool,
  body: (match (CalculusCheck::listAt [Bool] flags index)
    | Option.some flag => flag
    | Option.none => false),
  executable: true
}

definition{
  name: callsFail,
  parameters: [calls : (List Nat), flags : (List Bool)],
  result: Bool,
  body: (match calls
    | List.nil => false
    | List.cons head tail => (if (flagAt flags head) then true else (callsFail tail flags))),
  executable: true,
  recursive_argument: calls
}

definition{
  name: closuresFail,
  parameters: [closures : Clos, fnTy : CalculusCheck::Ty, flags : (List Bool)],
  result: Bool,
  body: (match closures
    | Clos.ClNil => false
    | Clos.ClCons known member captured rest => (if (CalculusCheck::tyEq known fnTy) then (if (flagAt flags member) then true else (closuresFail rest fnTy flags)) else (closuresFail rest fnTy flags))),
  executable: true,
  recursive_argument: closures
}

definition{
  name: appliesFail,
  parameters: [applies : CalculusCheck::Ty, closures : Clos, flags : (List Bool)],
  result: Bool,
  body: (match applies
    | CalculusCheck::Ty.TUnit => false
    | CalculusCheck::Ty.TBool => false
    | CalculusCheck::Ty.TNat => false
    | CalculusCheck::Ty.TInt => false
    | CalculusCheck::Ty.TFixed d0 => false
    | CalculusCheck::Ty.TString => false
    | CalculusCheck::Ty.TBytes => false
    | CalculusCheck::Ty.TOrdering => false
    | CalculusCheck::Ty.TOption d0 => false
    | CalculusCheck::Ty.TResult d0 d1 => false
    | CalculusCheck::Ty.TList d0 => false
    | CalculusCheck::Ty.TPair d0 d1 => false
    | CalculusCheck::Ty.TAdt d0 => false
    | CalculusCheck::Ty.TFn d0 d1 => false
    | CalculusCheck::Ty.TNil => false
    | CalculusCheck::Ty.TCons head tail => (if (closuresFail closures head flags) then true else (appliesFail tail closures flags))),
  executable: true,
  recursive_argument: applies
}

definition{
  name: failStep,
  parameters: [facts : Facts, closures : Clos, flags : (List Bool), index : Nat],
  result: (List Bool),
  body: (match facts
    | Facts.FactsNil => (nil Bool)
    | Facts.FactsCons own calls applies rest => (cons (if (if own then true else (if (flagAt flags index) then true else (callsFail calls flags))) then true else (appliesFail applies closures flags)) (failStep rest closures flags (index + 1)))),
  executable: true,
  recursive_argument: facts
}

definition{
  name: failNext,
  parameters: [facts : Facts, closures : Clos, flags : (List Bool)],
  result: (Option (List Bool)),
  body: (let next : (List Bool) := (failStep facts closures flags 0); (if (boolsEq next flags) then (ctor Option.none [(List Bool)]) else (ctor Option.some [(List Bool)] next))),
  executable: true
}

definition{
  name: deadFunctions,
  parameters: [
    spec : Bytes,
    adts : CalculusCheck::Ty,
    parameters : (List (List Nat)),
    types : CalculusCheck::Ty,
    results : CalculusCheck::Ty,
    bodies : CalculusCheck::Expr,
    sites : CalculusCheck::Ty
  ],
  result: (List Bool),
  body: (let program : CalculusCheck::Program := (record CalculusCheck::Program {spec := spec, adts := adts, parameters := parameters, types := types, results := results, bodies := bodies}); (let count : Nat := (CalculusCheck::exprChainLength bodies); (let registered : CalculusCheck::Ty := (register (register (register (ctor CalculusCheck::Ty.TNil) types) results) adts); (let initial : Collect := (collectFunctions program record{
    type: Collect,
    fields: [{field: known, value: registered}, {field: closures, value: (ctor Clos.ClNil)}],
    type_arguments: []
  } bodies 0 (nil Bool) sites); (let known : CalculusCheck::Ty := (project initial known); (let owners : CalculusCheck::Ty := (tyAppend (adtOwners (CalculusCheck::tyChainLength adts)) known); (let direct : Edges := (allEdges program known (project initial closures) owners (ctor Edges.EdgeNil)); (let edges : Edges := (closeEdges direct (1 + (prim multiply (CalculusCheck::tyChainLength owners) (CalculusCheck::tyChainLength owners) : Nat))); (let live : Liveness := (first (prim iterate_until (fun [bodies, edges, known, owners, program, sites] (state : Liveness) => (livenessStep program bodies known edges owners sites state)) (1 + count) record{
    type: Liveness,
    fields: [
      {field: dead, value: (falses count)},
      {field: closures, value: (project initial closures)}
    ],
    type_arguments: []
  } : (Liveness * Bool))); (project live dead)))))))))),
  executable: true,
  production: {effects: [allocation, overflow, recursion], targets: [rust-std]}
}

definition{
  name: fallibleFunctions,
  parameters: [
    spec : Bytes,
    adts : CalculusCheck::Ty,
    parameters : (List (List Nat)),
    types : CalculusCheck::Ty,
    results : CalculusCheck::Ty,
    bodies : CalculusCheck::Expr,
    sites : CalculusCheck::Ty,
    dead : (List Bool)
  ],
  result: (List Bool),
  body: (let program : CalculusCheck::Program := (record CalculusCheck::Program {spec := spec, adts := adts, parameters := parameters, types := types, results := results, bodies := bodies}); (let count : Nat := (CalculusCheck::exprChainLength bodies); (let registered : CalculusCheck::Ty := (register (register (register (ctor CalculusCheck::Ty.TNil) types) results) adts); (let initial : Collect := (collectFunctions program record{
    type: Collect,
    fields: [{field: known, value: registered}, {field: closures, value: (ctor Clos.ClNil)}],
    type_arguments: []
  } bodies 0 dead sites); (let facts : Facts := (reachFunctions program bodies 0 sites); (first (prim iterate_until (fun [facts, initial] (flags : (List Bool)) => (failNext facts (project initial closures) flags)) (1 + count) (falses count) : ((List Bool) * Bool)))))))),
  executable: true,
  production: {effects: [allocation, overflow, recursion], targets: [rust-std]}
}
\end{semanticmodule}
\end{lexlean}
